\documentclass{article}
\usepackage{amssymb, amsmath}
\usepackage[left=1.5cm, right=2cm]{geometry}
\begin{document}

\textbf{Hash Functions}\\
Input $m$: binary string\\
Output $h(m)$:
\begin{itemize}
	\item "Short" (n-bit) binary string (aka \textbf{message digest} or hash)
\end{itemize}
Efficiently computable\\
Applications: cryptography, security, efficiency\\
\underline{Keyed} $h_k(m)$, where the key is public, or \underline{unkeyed} $h(m)$\\

\textbf{Motivation: Hashing for Efficiency}\\
Input: large set (e.g. ints or strs)\\
Goal: map "randomly" to few bins\\

\textbf{Security Goal: Collision Resistance}:\\
A \textbf{collision:} two inputs (e.g. names) with the same hash: $h('Bob')=h('Phil')$\\
Every hash has collissions, since $|$input$| >> |$output$|$\\
Collision resistance: hard to \textbf{find} these collisions
\begin{itemize}
	\item Note: attacker can always try names randomly until a collision is found
	\item But this should be ineffective: must try about (on average) $N$ names (number of bins)
\end{itemize}

\textbf{Collision Resistant Hash Function (CRHF)}\\
$h$ is CRHF if it is hard to find collisions $h(x)=h(x')$ for $x\neq x'$
\begin{itemize}
	\item Note: attacker can always try names randomly until collision is found
	\item But this should be ineffective: must try about $|Range|$ values
\end{itemize}
\textbf{Definition} (Keyless Collision Resistant Hash Function (CRHF)). A keyless hash function $h^n(.) : \{0,1\}^* \rightarrow \{0,1\}^n$ is \textit{collision-resistant} if for every efficient (PPT) algorithm $A$, the advantage $\varepsilon^{CRHF}_{h,A}(n)$ is negligible in $n$, i.e. smaller than any positive polynomial for sufficiently large $n$ (as $n\rightarrow \infty$), where: \[
	\varepsilon^{CRHF}_{h,A}(n)=Pr\left[(x,x')\leftarrow A(1^n) s.t. (x\neq x') \wedge (h^n(x)=h^n(x'))\right]
\]

\textbf{Keyless CRHF Do Not Exist}\\
$|Range| << |Domain|$ so there \underline{is} a collision where $h(x')=h(x), x\neq x'$\\
Keys are powerful because they allow you to change the terms on your adversary.\\
For a keyless CRHF there \underline{is} a PPT algorithm A that can always output a collision: $A(1^n)=\{return\ x,x'\}$\begin{itemize}
	\item Proof: intuitively, since the function is fixed (same input-output mapping), a collision instance can be hardcoded int he attacker algorithm and just output that collision and win the security game.
\end{itemize}

\textbf{Keyed CRHF}\\
$k\overset{\$}{\leftarrow}\{0,1\}^n \leftarrow \mbox{Adversary} \leftarrow$ $(x,x')$ If $x\neq x'$ and $h(x)=h(x')$\\

\textbf{Generic Collision Attacks}\\
An attacker that runs in exponential time can always find a collision (i.e. non-PPT attacker)
\begin{itemize}
	\item Easy: find collisions in $2^n$ time by trying $2^n+1$ distinct inputs (compute their hash and locate a collision)
\end{itemize}
An attacker finds a collision with $2^{-n}$ probability (negligible probability).
\begin{itemize}
	\item Choose $x$ and $x'$ at random and check if they produce a collision.
\end{itemize}

\textbf{The Birthday Paradox}\\
The birthday paradox states that the expected number $q$ of hashes (i.e. hash queries) until a collision is found in $O(2^{\frac{n}{2}})$ not $O(2^n)$. It is $q\lessapprox2^{n/2}*\sqrt{\frac{\pi}{2}}\lessapprox1.254*2^{n/2}$\\
For 80 bits of effective security, use $n=160$
\begin{itemize}
	\item So to defend against an attacker who can perform $2^{80}$ hashes set the digest length to at least 160 bits (so the range has a size of $2^{160}$ digests)
\end{itemize}

\textbf{The Birthday Attack}\\
Probability of NO birthday-collision:\begin{itemize}
	\item Two people: (364/365)
	\item Three people (364/365)(364/365)
	\item ...
	\item $n$ people: $\prod^{n-1}_{i=1}\frac{365-i}{365}$
	\item So for 23 individuals, the probability of a birthday collision is about .5
\end{itemize}

\textbf{Collision-Resistance: Applications}\\
Integrity (of object/file/message)\begin{itemize}
	\item Send $hash(m)$ securely to validate $m$
	\item Later, we will see how a hash funciton can be used to construct a MAC (called HMAC)
\end{itemize}
Hash-then-Sign\begin{itemize}
	\item Instead of signing $m$, sign $hash(m)$ (more efficient!)
\end{itemize}
Blockchains\\

\textbf{Other Notions of Security}\\
Collision resistance provides the strngest guarantee\begin{itemize}
	\item Gives mroe freedom to the adversary; the adversary always wins if it finds any two inputs with the same digest.
	\item No conditions on these two inputs other than being in the domain of the hash function.
\end{itemize}
\textbf{Hash Function}: is a mapping that takes an input (VIL) and maps it to a fixed length $n$ string where $n$ is the security parameter.\\
Other security notions (but sufficient for many applications):\begin{itemize}
	\item Target collision resistance (TCR)
	\item Second preimage resistance
	\item First preimage resistance
\end{itemize}
Birthday paradox (or attack) does not work against these weaker notions. It is for collision resistance; find \textit{any} two inputs that collide!\\

\textbf{TCR and Birthday Paradox?}\\
TCR: Adversary has to select the target \textit{before} knowing the key\\

\textbf{We (mostly) focus on keyless hash}\\
Although there are no CRHFs and theory papers focus on keyed hash\\
But...\begin{itemize}
	\item It's a bit less complicated and easier to work with.
	\item No need to consider both ACR (Any Collision Resistance, which is the same as CRHF) and TCR. Why?
	\item Modifying to CRHF is quite trivial: just make it keyed
	\item Usually used in practice: code libraries, standards, ...
\end{itemize}

\textbf{$2^{nd}$-Preimage-Resistant Hash (SPR)}\\
\underline{$y = h(x)$; image = preimage}\\
Hard to find collision with a \underline{specific random $x$}\\
So, the attacker's advantage (probability of finding such a collision) in the following game is negligible.\\
Here, adversary has no freedom over the first preimage $x$, whereas in TCR, the attacker can choose $x$ (and then the second x is after knowing the hash function)\\
Example: hashing software from developer is SPR, since the software hash is like the first preimage. The attacker must find a collision against this particular message (e.g. the software)\\

\textbf{CRHF/SPR vs. Applications}\\
CRHF secure for signing, SW-distribution\\
How about SPR hash (weak CRHF)?\begin{itemize}
	\item SW-distribution? YES
	\item Hash-then-sign? NO
\end{itemize}
Why?\\
Attacker can't impact SW to be distributed, but the attacker may be able to impact the signed message!\\
Since h is not necessarily CRHF, its signature can be forged (the tag for two collisions is the same)\\

\textbf{SPR: Collisions to Chosen Messages}\\
Or: Alice and Mal, the corrupt lawyer\\
Mal finds two "colliding wills," GoodW and BadW:\begin{itemize}
	\item GoodW: contents agreeable to Alice
	\item $h(GoodW)=h(BadW)$
	\item Alice signs good will: $Sign(h(GoodW))$
	\item $Sign(h(GoodW))=Sign(h(BadW))$
\end{itemize}
Is such an attack realistic? Or is SPR enough "in practice?" Yes.\\
We usually rely on having a specific format for messages, like a fixed prefix chosen by the attacker.\\

\textbf{Examples}
\begin{enumerate}
	\item Let $h_k$ be a keyed CRHF. Is $h'_k=h_k(h_k(x))$ a CRHF? Why?
	\item For $x$ parsed as $x=x_1||x_2||x_3$, let $h(x_1||x_2||x_3)=x_1+x_2+x_3 \% p$, is $h$ a CRHF? Why? Is it an SPR? Why?
	\item Let $h_k(m)$ be a TCR function. Construct $h_k'(m)=0^n$ if $m[1:|k|]=k$ and $h_k(m)$ otherwise (recall that the hash function input can be of any length). Is $h_k'$ a CRHF? Why? Is $h'_k$ a TCR? Why?
		\begin{itemize}
			\item Not a CRHF: A knows $k$ before outputting $m$ and $m'$, A can pick any two messages where the first $|k|$ bits are the key.
			\item TCR: Attacker doesn't know key beforehand. $m$ chosen in advance must start with the key for the attacker to be successful, but the chances of this happening are $\frac{1}{2^{|k|}}$, which is negligible.
		\end{itemize}
\end{enumerate}

\textbf{One-Way Function (OWF/First Preimage Resistance)}\\
One-way function or first preimage resistance: given $h(k)$ for a \underline{random} $x$, it is hard to find $x$, \underline{or any $x'$} s.t. $h(x')=h(x)$ (like inverting the hash func)\\ 
Given the output, cannot find the input (or a collision).\\
Compare to:\begin{itemize}
	\item Collision-Resistance (CRHF): hard to find any collision, i.e. any $(x,x')$ s.t. $h(x')=h(x), x\neq x'$
	\item Second-Preimage Resistance (SPR): hard to find a collision with \underline{random} $x$, i.e. $x'$ s.t. $h(x')=h(x), x\neq x'$
\end{itemize}

\textbf{Application: One-time Password Authentication}\\
Used with OWF, since it can't be inverted to glean passwd\\
One-time password auth:\begin{itemize}
	\item Select random $x$ : "one-time password" (keep secret)
	\item Validate using non-secret "one-time validation token:" $h(x)$
\end{itemize}

\textbf{Randomness Extraction}\\
It is a desired property of hash funcs.\\
"If input is sufficiently random, then the output is random"\\
\underline{Randomness extraction}: if any $m$ input bits are random $\rightarrow$ all $n$ output bits are pseudorandom, for sufficiently large $m$\\
ROM (Random Oracle Model): assume hash function behaves as a true random function (difference from randomness extraction, which is not an assumption)\\

\textbf{Exercise}\\
Let $h_1,h_2$ be \underline{both} CRHF and OWF\\
Use them to construct $h_{CRHF}$: CRHF but not OWF and $h_{OWF}$ but not CRHF\\
One possible solution:\begin{itemize}
	\item \[h_{CRHF}(m)=\begin{cases}
		m & if\ |m| = n \\
		h_1(m) & otherwise
	\end{cases}\]\\
	Since adversary selects the length of the message, it can set it to length $n$ which forces it to branch to 1 every time in the security game.
	\item \[h_{OWF}(m)=\begin{cases}
		h_1(m) & if\ |m|=n \\
		h_1(m_{1..n}\oplus h_2(m')) & if\ m=m_{1..n}||m'
	\end{cases}\]\\
	Not OWF. Once an attacker computes an arbitrary $h_2(m')$, they can find an $m_{1..n}$ that produces the same output when $\oplus$'d with $h_2(m')$.
\end{itemize}

\textbf{Security Notions}\\
CRHF $>$  TCR $>$ SPR\\
OWF is separate to all of these.\\


\end{document}
